\scanpage{267}
\mspage{79}
% archive sheet number: 267
Pour $f$ fixé, on va construire une catégorie auxiliaire $K(f)$, de telle façon que pour \unsure{tte|toute} catégorie $\mathcal M$, les systèmes $(G,F,u)$ d'un $\mathcal M$ préfaisceau \inserted{$G$} sur $J$, d'un autre $F$ sur $I$, et d'\unsure{un|une} $f$-\unsure{hom|flèche} $G\unsure{\text{---}|\longrightarrow}F$, s'identifient (\unsure{fonctoriellement|fonctoriell\textsuperscript{t}}) aux $\mathcal M$-préfaisceaux sur $K(f)$.

Ceci dit, on va essayer d'interpréter les flèches \inserted{(1)} (squelettiques, mais peu importe ici --- tous les \inserted{$\mathbb D$-}squelettes sont \inserted{maintenant} remplumés) en termes de la géométrie de $K(f)$

\[
\left\{\begin{aligned}
&\mathrm{Ob}\,K(f)=\mathrm{Ob}\,K(I)\amalg\mathrm{Ob}\,K(J)\\
&I,\ J\ \text{s'identifient à des sous-catégories pleines de}\\
&\qquad K(f),\ \text{il n'y a pas de flèches d'un }j\in J\text{ en un}\\
&\qquad i\in I.\quad(\text{Donc }I\text{ est ouverte dans }K(f),\ J\ \deleted{\unsure{f}}\ \text{\unsure{son}}\\
&\qquad\text{complémentaire fermé})\\
&\text{Pour }i\in\mathrm{Ob}\,I,\ j\in\mathrm{Ob}\,J,\ \text{on pose}\\
&\qquad\operatorname{Hom}_{\deleted{\unsure{K_J|J}}K(f)}(i,j)=\operatorname{Hom}_J(f(i),j)
\end{aligned}\right.
\]

La composition des flèches dans $K(f)$ est évidente. On va définir une rétraction canonique

\[
\bar f:K(f)\longrightarrow J
\]

de $K(f)$ sur sa sous-catégorie fermée $J$. Donc $\bar f$ est connue sur $\mathrm{Ob}\,J$, et sur $\mathrm{Ob}\,I$ ce sera $f$ :

\[
\left\{\begin{aligned}
\bar f(j)&=j&&\text{si }j\in\mathrm{Ob}\,J\\
\bar f(i)&=f(i)&&\text{si }i\in\mathrm{Ob}\,I
\end{aligned}\right.
\]

\deleted{Pour les flèches} Il reste à définir, pour $x,y\in\mathrm{Ob}\,K(f)$
